 \chapter{Opérations}\label{ch:fm:operations}

Le 28 mai 2024, la plupart des transactions sur titres américains sont passées d'un règlement-livraison à deux jours ouvrés après la transaction à un règlement-livraison à un jour. Une société qui disposait d'une journée pour trouver une
instruction de règlement manquante avait désormais la soirée du jour de la transaction~: une nouvelle règle exigeait que les allocations, les confirmations et les affirmations soient
achevées avant la fin du jour de la transaction, et l'heure limite d'affirmation du secteur était fixée à 21 h, heure de l'Est. Le rapport sectoriel publié ensuite constatait
que près de 95~\% des transactions respectaient cette heure limite, contre 73~\% en janvier 2024, et que les taux de suspens étaient restés au niveau qu'ils avaient sous le
cycle à deux jours. Le travail n'avait pas disparu. Il s'était concentré sur moins d'heures, et chaque équipe des opérations qui ne voulait pas plus de suspens devait être dimensionnée pour la
fenêtre plus courte.

\section{Front, middle et back office}

Une transaction est une promesse d'échanger des titres contre espèces à une date ultérieure. Entre la promesse et l'échange s'intercalent les étapes du chapitre sur le règlement-livraison
du livre~1 et de la plateforme post-marché du livre~15~: saisie des transactions (trade capture), enrichissement, confirmation, compensation, instruction de règlement, règlement-livraison et le
rapprochement (reconciliation) de chaque enregistrement avec celui de chaque contrepartie. Une société de trading répartit le travail en trois groupes.

\begin{definition}[Front, middle et back office]\label{def:fm:operations:offices}
Le \emph{front office}\index{front office} prend et exécute les décisions de trading~: les traders, les chercheurs et les commerciaux. Le \emph{middle
office}\index{middle office} contrôle et rend compte de ce que le front office a fait, indépendamment de lui~: validation des transactions, reporting du P\&L et des
positions, contrôle des risques et le product control. Le \emph{back office}\index{back office} mène les transactions jusqu'à leur achèvement~: confirmation,
compensation, règlement-livraison, mouvements d'espèces et de collatéral, opérations sur titres et rapprochement.
\end{definition}

La division est aussi un contrôle. La séparation des tâches (livre~6) exige que la personne qui passe une transaction ne l'enregistre, ne la confirme ni ne la règle,
et que la personne qui rend compte de son P\&L ne la négocie pas~: le \omterm{def:fm:operations:offices}{middle office} existe pour que les chiffres du \omterm{def:fm:operations:offices}{front office} soient vérifiés par quelqu'un qui
n'est pas rémunéré sur ces chiffres. Le risque opérationnel (livre~6) se loge dans les passages de relais entre les trois.

\begin{omfigure}
\begin{tikzpicture}[font=\footnotesize,>=stealth,box/.style={draw,rounded corners,align=center,minimum height=0.75cm,text width=2.2cm}]
\node[box,fill=omDef!15] (fo) at (0,0) {front office\\{\scriptsize exécuter}};
\node[box,fill=omProp!15] (mo) at (3.1,0) {middle office\\{\scriptsize valider, reporter}};
\node[box,fill=omThm!15] (bo) at (6.2,0) {back office\\{\scriptsize confirmer, régler}};
\node[box,fill=gray!15] (st) at (9.3,0) {règlement-livraison\\{\scriptsize titres contre espèces}};
\draw[->] (fo) -- (mo); \draw[->] (mo) -- (bo); \draw[->] (bo) -- (st);
\draw[thick] (0,-1.3) -- (9.3,-1.3);
\foreach \x/\l in {0/T,4.65/{T+1},9.3/{T+2}} {\draw (\x,-1.2) -- (\x,-1.4) node[below] {\l};}
\draw[<->,omThm,thick] (0,-2.2) -- node[below,font=\scriptsize] {cycle à deux jours~: corriger les écarts pendant T+1} (9.3,-2.2);
\draw[<->,omDef,thick] (0,-2.9) -- node[below,font=\scriptsize] {cycle à un jour~: affirmer d'ici le soir de T} (4.65,-2.9);
\end{tikzpicture}

\omcaption{Les trois fonctions le long de la vie d'une transaction, et la fenêtre dans laquelle un écart doit être corrigé avant le règlement-livraison, sous un cycle à deux jours et un cycle à
un jour. Schéma de principe.}\label{fig:fm:operations:offices}
\end{omfigure}

\begin{dated}{2026-09}{Cycles de règlement-livraison}\label{dat:fm:operations:cycle}
\textbf{États-Unis}~: la SEC, dans le Release 34-96930 (février 2023), a modifié la Rule 15c6-1 pour qu'un broker-dealer ne puisse pas, pour la plupart des titres, convenir
de régler plus tard que le premier jour ouvré après la transaction, sauf si les parties en conviennent expressément autrement au moment de la transaction~; la mise en conformité a commencé
le 28 mai 2024. La nouvelle Rule 15c6-2 exige que les allocations, les confirmations et les affirmations soient achevées dès que cela est techniquement possible et au
plus tard à la fin du jour de la transaction. Le SIFMA, l'ICI et la DTCC ont indiqué en septembre 2024 que près de 95~\% des transactions respectaient l'heure limite d'affirmation
de 21 h, heure de l'Est (73~\% en janvier 2024), et que les taux de suspens de juillet 2024, de 2.12~\% dans le système continuous net settlement et de 3.31~\% en dehors, étaient conformes à ceux du T+2.
\textbf{Union européenne}~: la Regulation (EU) 2025/2075 fait passer le cycle de règlement-livraison de la Central Securities
Depositories Regulation, de deux jours ouvrés à un jour à compter du 11 octobre 2027.
\end{dated}

\section{Rapprochement et écarts}

Les enregistrements d'une société relatifs à une transaction, à une position ou à un solde en espèces doivent concorder avec ceux de sa contrepartie, de sa chambre de compensation, de son \omterm{def:fm:the-asset-manager-and-the-fund:admin}{dépositaire} et de ses
autres systèmes internes. Le rapprochement les compare~; une différence est un écart de rapprochement (break, livre~15), et son ancienneté depuis sa découverte est le vieillissement des écarts
que l'équipe des opérations publie chaque jour. Les écarts ont des causes~: un prix ou une quantité erronés à la saisie, une instruction de règlement absente ou périmée, une
allocation non transmise, une opération sur titres enregistrée autrement. La plupart se résolvent en quelques minutes~; quelques-uns prennent des jours~; un écart encore ouvert lorsque les titres
doivent être transférés devient un suspens de règlement-livraison.

L'équipe qui les résout est une file d'attente (livre~4)~: les écarts arrivent, attendent qu'une personne soit libre et sont traités. Supposons que les écarts arrivent au taux $\lambda$ par heure,
que chacun prenne un temps exponentiel de moyenne $1/\mu$ pour être résolu, et que $c$ personnes les traitent, premier arrivé, premier servi. La file est stable si
$\lambda<c\mu$. Un écart est en retard si son temps dans le système, attente comprise plus traitement, dépasse le délai $D$ restant avant le règlement-livraison.

\begin{proposition}[Écarts en retard dans une équipe de $c$]\label{prop:fm:operations:late}
Dans la file M/M/c stationnaire avec $a=c\mu-\lambda>0$, notons $C$ la probabilité qu'un écart arrivant doive attendre (formule d'Erlang C). La
probabilité qu'un écart ne soit pas résolu dans le délai $D$ est
\[
P(T>D)=(1-C)\,e^{-\mu D}+C\,\frac{a\,e^{-\mu D}-\mu\,e^{-aD}}{a-\mu}\qquad(a\neq\mu),
\]
et $P(T>D)\ge e^{-\mu D}$ pour tout $c$~: si $e^{-\mu D}>1-s$, aucune équipe n'est assez nombreuse pour résoudre à temps une part $s$ des écarts.
\end{proposition}

\begin{proof}[Preuve partielle]
Dans la file FCFS M/M/c, un écart arrivant attend avec probabilité $C$ et, sachant qu'il attend, son attente suit une loi exponentielle de paramètre $a$ (un
résultat classique, \admitted). Son temps de traitement $S$ suit une loi exponentielle de paramètre $\mu$, indépendante de l'attente. Sans attente, $P(S>D)=e^{-\mu D}$~; avec une attente $W$,
$P(W+S>D)$ est la queue de la somme de deux exponentielles indépendantes, qui est la fraction donnée. Les deux termes valent au moins $e^{-\mu D}$, puisque
$W\ge0$.
\end{proof}

\omcode{code/firm/opsmetrics/firm_opsmetrics.py}{39}{68}{Probabilité d'attente tirée de la loi stationnaire de firm.queues, part des écarts en retard de
la \cref{prop:fm:operations:late}, et la plus petite équipe pour un niveau de service.}\label{lst:fm:operations:late}

Le plancher est le point essentiel de la proposition. Le personnel réduit l'attente, pas le travail~: dès qu'aucun écart n'attend, la part des écarts en retard est celle des écarts dont
le seul traitement dépasse la fenêtre. Sous un cycle à deux jours, la fenêtre est longue et le plancher invisible~; sous un cycle à un jour, il peut être la
contrainte active.

\section{Tutoriel~: un jour de moins}\label{tut:fm:operations:tut}

\begin{tutorial}
\textbf{Objectif.} Dimensionner une équipe de traitement des écarts pour un cycle à deux jours et pour un cycle à un jour, et chiffrer les suspens qu'elle laisse passer. \textbf{État final~:}
la \cref{fig:fm:operations:late}, la \cref{fig:fm:operations:cost} et l'effectif nécessaire pour un niveau de service de 95~\%.
\begin{tutsteps}
\item \textbf{Le flux.} 120 écarts par jour arrivent sur une journée de huit heures ($\lambda=15$ par heure) et demandent en moyenne 40 minutes de travail ($\mu=1.5$)~;
la charge offerte est de $\lambda/\mu=10$ personnes (à titre illustratif).
\item \textbf{Les fenêtres.} Un écart doit être résolu dans un délai de 10 heures ouvrées sous T+2 et de 2.5 heures sous T+1, avant l'heure limite du soir.
\item \textbf{Les suspens.} La moitié des écarts en retard devient des suspens, portant sur des transactions de 5~millions de dollars, restant en suspens pendant deux jours moyennant une pénalité de 1 point de base par jour
(le taux de l'UE pour les actions liquides, une donnée d'entrée) plus 1\,500~\$ de coût interne~: 2\,500~\$ par suspens. Une personne coûte 120\,000~\$ par an.
\item \textbf{Le dimensionnement.} \texttt{fm\_ops.staff95(cycle)} applique \texttt{firm.opsmetrics.staff\_for}~; \texttt{fm\_ops.table(cycle)} chiffre
les équipes selon leur effectif et le coût des suspens.
\end{tutsteps}
\end{tutorial}

\begin{omfigure}
\begin{tikzpicture}
\begin{axis}[width=10.5cm,height=5.2cm,xlabel={\small personnes dans l'équipe},ylabel={\small écarts en retard (\%)},xmin=11,xmax=20,ymin=0,ymax=10,
  tick label style={font=\footnotesize},grid=major,grid style={gray!20},table/col sep=comma,
  legend cell align=left,legend style={font=\scriptsize,at={(0.5,-0.3)},anchor=north,/tikz/every even column/.append style={column sep=8pt}},legend columns=3]
\addplot[omThm,thick,mark=*,mark size=1.2pt] table[x=team,y=d10]{figdata/desk/13-operations/late.csv};
\addplot[omProp,thick,dashed,mark=square*,mark size=1.2pt] table[x=team,y=d4]{figdata/desk/13-operations/late.csv};
\addplot[omDef,thick,mark=triangle*,mark size=1.5pt] table[x=team,y=d25]{figdata/desk/13-operations/late.csv};
\draw[gray,dotted] (axis cs:11,5) -- (axis cs:20,5);
\draw[omDef,densely dotted] (axis cs:11,2.35) -- (axis cs:20,2.35);
\legend{fenêtre de 10 heures (T+2),fenêtre de 4 heures,fenêtre de 2.5 heures (T+1)}
\end{axis}
\end{tikzpicture}

\omcaption{La part des écarts non résolus dans la fenêtre de règlement, selon la taille de l'équipe, pour 120 écarts par jour de 40 minutes de travail. Avec 2.5 heures, aucune
équipe ne descend sous le plancher $e^{-\mu D}=2.35~\%$ (pointillé inférieur) fixé par le travail lui-même~; le pointillé supérieur est le 5~\% d'un niveau de service de 95~\%.
Données~: \texttt{firm.opsmetrics.p\_late}.}\label{fig:fm:operations:late}
\end{omfigure}

Sous T+2, onze personnes, la plus petite équipe stable, résolvent tous les écarts sauf quelques-uns sur un million dans la fenêtre~: la part en retard est de 0.0003~\% et
le coût annuel des suspens est nul. Sous T+1, onze personnes laissent 8.37~\% des écarts en retard, soit 1\,255 suspens par an pour un coût de 3.14~millions de dollars~; douze sont nécessaires
pour le niveau de service de 95~\% (3.38~\% en retard, 1.27~million de dollars) et treize minimisent à la fois les effectifs et le coût des suspens (2.69~\% en retard, 2.57~millions de dollars au total,
contre 1.32~million de dollars sous T+2). Au-delà de treize, les personnes n'achètent presque plus rien~: la part en retard s'approche du plancher de 2.35~\%
(\cref{fig:fm:operations:cost}). L'attente est la part la moins importante du problème~: avec treize personnes, 28.5~\% des écarts attendent au moins un instant.

\begin{omfigure}
\begin{tikzpicture}
\begin{axis}[width=10.5cm,height=5.2cm,xlabel={\small personnes dans l'équipe},ylabel={\small millions de dollars par an},xmin=11,xmax=18,ymin=0,ymax=5,
  tick label style={font=\footnotesize},grid=major,grid style={gray!20},table/col sep=comma,
  legend cell align=left,legend style={font=\scriptsize,at={(0.5,-0.3)},anchor=north,/tikz/every even column/.append style={column sep=8pt}},legend columns=3]
\addplot[omThm,thick,dashed] table[x=team,y=staff]{figdata/desk/13-operations/cost.csv};
\addplot[omProp,thick,dotted] table[x=team,y=fail]{figdata/desk/13-operations/cost.csv};
\addplot[omDef,very thick,mark=*,mark size=1.2pt] table[x=team,y=total]{figdata/desk/13-operations/cost.csv};
\legend{coût des effectifs,coût des suspens,total}
\end{axis}
\end{tikzpicture}

\omcaption{Coût des effectifs, coût des suspens et leur somme avec un cycle d'un jour. Le total est le plus bas à treize personnes, soit 2.57~millions de dollars par an, où le coût des suspens reste de
1.01~million de dollars~: le plancher de la \cref{prop:fm:operations:late} ne peut pas être levé par l'effectif. Données~: \texttt{fm\_ops.table}.}\label{fig:fm:operations:cost}
\end{omfigure}

\begin{table}[ht]
\centering\small
\begin{tabular}{>{\raggedright\arraybackslash}p{2.0cm}>{\centering\arraybackslash}p{2.3cm}>{\centering\arraybackslash}p{2.3cm}>{\centering\arraybackslash}p{2.3cm}>{\centering\arraybackslash}p{2.3cm}>{\centering\arraybackslash}p{2.3cm}}
\toprule
scénario & équipe pour 95~\% & équipe la moins coûteuse & coût des effectifs & coût des suspens & total \\
\midrule
T+2 & 11 & 11 & 1.32 & 0.00 & 1.32 \\
T+1 & 12 & 13 & 1.56 & 1.01 & 2.57 \\
T+1, moitié des écarts & 7 & 7 & 0.84 & 0.58 & 1.42 \\
T+1, 30 minutes de travail & 9 & 10 & 1.20 & 0.30 & 1.50 \\
\bottomrule
\end{tabular}
\caption{Dimensionnement d'une équipe de traitement des écarts sous chaque cycle, et deux façons de modifier le travail lui-même ($\$$ millions par an). Données~:
\texttt{fm\_ops.scenarios}.}\label{tab:fm:operations:scenarios}
\end{table}

Le remède consiste à modifier le travail, non l'équipe (\cref{tab:fm:operations:scenarios}). Réduire de moitié le nombre d'écarts, en automatisant les étapes qui les provoquent
(instructions de règlement permanentes, allocation et appariement électroniques à la date de transaction), demande sept personnes et coûte 1.42~million de dollars~; rendre
le travail plus rapide, à 30 minutes par écart, en demande dix et coûte 1.50~million de dollars. Les deux options annulent presque le coût du cycle plus court, et l'effectif seul
ne suffit pas.

\begin{remark}[Ce que la file d'attente laisse de côté]\label{rem:fm:operations:batch}
Les écarts n'arrivent pas uniformément~: beaucoup apparaissent ensemble, au moment où les rapprochements de fin de journée sont exécutés, ce qui est le pire moment sous un cycle d'un jour. Une file
d'attente à arrivées groupées en clôture donne des parts en retard plus grandes que le flux régulier du chapitre, et plaide pour déplacer le rapprochement en temps réel
(livre~15) plutôt que pour davantage de personnel le soir.
\end{remark}

\begin{omfigure}
\begin{tikzpicture}
\begin{axis}[width=10.5cm,height=5cm,ybar,bar width=8pt,ylabel={\small part des écarts (\%)},ymin=0,ymax=60,xtick={0,1,2,3,4},
  xticklabels={{$<$30 min},{30--60 min},{1--2 h},{2--4 h},{$>$4 h}},xmin=-0.6,xmax=4.6,
  tick label style={font=\footnotesize},grid=major,grid style={gray!20},table/col sep=comma,
  legend cell align=left,legend style={font=\scriptsize,at={(0.5,-0.25)},anchor=north,/tikz/every even column/.append style={column sep=8pt}},legend columns=3]
\addplot[fill=omDef!55,draw=omDef] table[x=pos,y=c11]{figdata/desk/13-operations/ageing.csv};
\addplot[fill=omProp!55,draw=omProp] table[x=pos,y=c12]{figdata/desk/13-operations/ageing.csv};
\addplot[fill=omThm!50,draw=omThm] table[x=pos,y=c14]{figdata/desk/13-operations/ageing.csv};
\legend{11 personnes,12 personnes,14 personnes}
\end{axis}
\end{tikzpicture}

\omcaption{Profil d'ancienneté des écarts résolus~: le temps écoulé entre l'arrivée et la résolution, pour trois tailles d'équipe. Davantage de personnes font passer les écarts
dans la première demi-heure~; la part de plus de quatre heures passe de 1.3~\% à 0.3~\%, sans toutefois atteindre zéro. Données~: \texttt{fm\_ops.ageing}.}\label{fig:fm:operations:ageing}
\end{omfigure}

\section{Opérations sur titres et événements de référence}

Un dividende, un split, une fusion, une augmentation de capital avec droits ou un coupon modifie la valeur d'une position ou sa nature. Le livre~1 décrit ces événements~; pour
les opérations, ce sont autant d'écarts en attente, puisque chaque système qui détient la position doit appliquer l'événement de la même manière, à la même date, et
la société doit trancher les événements à option (encaisser en espèces ou en titres, apporter ou non) avant une échéance fixée par l'agent de l'émetteur. Les contrôles sont
ceux du rapprochement, appliqués avant l'événement plutôt qu'après~: une source de référence (golden source) pour les données de référence, un calendrier des événements vérifié
par rapport aux positions concernées, et un responsable désigné pour chaque décision à option.

\section{Gestion des garanties}

Toute relation soumise à marge fait circuler le collatéral chaque jour~: la marge de variation avec les chambres de compensation (livre~1) et, sous une credit support annex (livre~2),
avec les contreparties bilatérales au-delà d'un seuil et pour des montants non inférieurs au montant minimal de transfert (livre~6).

\begin{definition}[Gestion des garanties, litige sur les garanties]\label{def:fm:operations:collateral}
La \emph{gestion des garanties}\index{gestion des garanties} est le processus quotidien de calcul, d'appel, de paiement, de réception, de valorisation et de substitution
du collatéral dû au titre des accords de marge d'une société, ainsi que de rapprochement de ses registres avec ceux de ses contreparties. Un \emph{litige sur les
garanties}\index{litige sur les garanties} survient lorsque les calculs des deux parties sur le montant dû diffèrent de plus d'une tolérance convenue~; le
montant non contesté est généralement payé pendant que l'écart est examiné.
\end{definition}

\begin{example}[Un appel contesté]\label{ex:fm:operations:call}
Une CSA prévoit un seuil de 10 millions de dollars et un montant minimal de transfert de 0.5 million de dollars~; la société détient 14.6 millions de dollars. La société valorise
l'exposition à 25.3 millions de dollars et appelle $25.3-10-14.6=0.7$ million de dollars. La contrepartie la valorise à 24.4 millions de dollars~: son chiffre représente une restitution de 0.2 million de dollars,
inférieure au montant minimal de transfert, de sorte qu'elle ne voit aucun appel. Les valorisations diffèrent de 3.6~\%, au-delà d'une tolérance de 2~\%~: l'appel est contesté, et la cause
est presque toujours une différence de population de transactions ou de valorisation de quelques transactions, c'est là que commence le rapprochement.
\end{example}

\omcode{code/firm/opsmetrics/firm_opsmetrics.py}{81}{106}{Un appel de collatéral sous le seuil et le montant minimal de transfert d'une CSA, le test de litige, et le feu tricolore
d'un indicateur clé de risque.}\label{lst:fm:operations:csa}

\section{Discipline en matière de règlement}

\begin{definition}[Régime de discipline en matière de règlement]\label{def:fm:operations:discipline}
Un \emph{régime de discipline en matière de règlement}\index{régime de discipline en matière de règlement} est l'ensemble des règles par lesquelles un marché dissuade les suspens de règlement-livraison et les traite~:
la déclaration des suspens, les pénalités en espèces imputées à la partie défaillante, et les buy-ins qui remplacent les titres non livrés.
\end{definition}

\begin{dated}{2026-09}{Le régime européen de discipline en matière de règlement}\label{dat:fm:operations:csdr}
Les règles de discipline en matière de règlement du Central Securities Depositories Regulation (Delegated Regulation (EU) 2018/1229) s'appliquent depuis le 1\up{er} février
2022. Les pénalités en espèces sont imputées par jour de suspens sur la valeur de l'instruction en défaut, à des taux fixés par le Delegated Regulation (EU) 2017/389~: 1.0
point de base pour les actions liquides, 0.5 pour les actions peu liquides, 0.25 pour les instruments des marchés de croissance des PME, 0.10 pour la dette souveraine, 0.20 pour les autres titres de dette,
0.15 pour la dette des marchés de croissance des PME et 0.5 pour les autres instruments~; un manque de liquidités est imputé au taux de la facilité de crédit à un jour de la banque centrale émettrice,
avec un plancher à zéro. Le Regulation (EU) 2023/2845 a fait des buy-ins obligatoires une mesure de dernier recours~: ils ne s'appliquent que si la Commission les introduit par
un acte d'exécution, après une analyse coûts-bénéfices menée par l'ESMA, lorsque les autres mesures n'ont pas réduit les suspens.
\end{dated}

Une pénalité de 1 point de base par jour sur une transaction de 5~millions de dollars représente 500~dollars par jour~: peu de chose au regard du P\&L d'un desk de trading, mais facturée sur chaque suspens et chaque jour,
et versée à la contrepartie qui n'a pas fait défaut. Les coûts plus importants d'un suspens sont ceux que le tutoriel chiffre comme internes~: le financement de la position non réglée,
les réclamations, le temps consacré et la dégradation de la relation. Un \omterm{def:fm:operations:kri}{indicateur clé de risque} permet de les rendre visibles avant qu'ils ne surviennent.

\begin{definition}[Indicateur clé de risque]\label{def:fm:operations:kri}
Un \emph{indicateur clé de risque}\index{indicateur clé de risque} (key risk indicator) est une métrique suivie par rapport à des seuils, parce qu'il évolue avant les pertes~: pour
les opérations, le nombre d'écarts de plus d'un jour, le taux de suspens, la part des transactions affirmées à la date de transaction, le nombre d'appels de marge contestés.
\end{definition}

\begin{method}[Un rapport quotidien des opérations]\label{met:fm:operations:report}
\begin{enumerate}
\item Choisir de cinq à dix \omterm{def:fm:operations:kri}{indicateurs clés de risque}, chacun assorti d'un seuil orange, d'un seuil rouge et d'un responsable.
\item Les calculer à partir des systèmes de référence, et non à partir des décomptes propres aux équipes.
\item Diffuser les feux tricolores chaque jour au \omterm{def:fm:operations:offices}{middle office} et au responsable des opérations, et chaque semaine, avec les tendances, au \omterm{def:fm:the-risk-management-function:cro}{comité des risques}.
\item Pour chaque feu rouge, consigner la cause et la correction~; pour chaque orange récurrent, modifier le processus, et non le seuil.
\end{enumerate}
\end{method}

Sur le rapport du chapitre (\texttt{fm\_ops.daily\_report}), quatorze écarts de plus d'un jour, un taux de suspens de 2.4~\% et 93~\% des transactions affirmées à la
date de transaction sont tous en orange, face à des seuils de 10, 2~\% et 95~\%~; un appel contesté est en vert.

\section{Construction~: métriques des opérations}\label{bld:fm:operations:opsmetrics}

\begin{build}
\bfield{Rôle} Les chiffres de l'équipe des opérations~: la file des écarts et son vieillissement, les effectifs pour un niveau de service, les coûts des suspens, les appels de collatéral et
litiges, et les \omterm{def:fm:operations:kri}{indicateurs clés de risque}.
\bfield{Interface} \texttt{firm.opsmetrics}~: \texttt{mmc\_stationary}, \texttt{p\_wait}, \texttt{erlang\_c}, \texttt{p\_late}, \texttt{staff\_for},
\texttt{ageing}, \texttt{fail\_cost}~; \texttt{Csa}, \texttt{call}, \texttt{disputed}~; \texttt{kri}, \texttt{report}.
\bfield{Règles} La file doit être stable~; un niveau de service inférieur au plancher de temps de travail est refusé~; les taux de pénalité sont des données, jamais des constantes~; un \omterm{def:fm:operations:kri}{indicateur clé
de risque} dont le rouge se situe en deçà de l'orange est lu dans le sens inverse.
\bfield{Tests d'acceptation} \texttt{code/firm/opsmetrics/tests/}~: la probabilité d'attente de la chaîne face à la formule d'Erlang~; le plancher~; le dimensionnement à
la frontière~; les appels, les litiges et les feux tricolores sur des chiffres saisis à la main.
\bfield{Pour aller plus loin} Arrivées par lots à la clôture~; types d'écarts ayant leurs propres temps de travail~; un circuit de collatéral avec substitution et vieillissement des
litiges.
\end{build}

\begin{omsources}
\item SEC Release 34-96930 (2023)~; 17 CFR 240.15c6-1 et 240.15c6-2.
\item SIFMA, ICI et DTCC, T+1 After Action Report, communiqué de presse du 12 septembre 2024.
\item Delegated Regulations (EU) 2017/389 et 2021/70~; Regulations (EU) 2023/2845 et 2025/2075.
\end{omsources}

\section{Exercices}

\begin{exercise}[$\star$]\label{exo:fm:operations:1}
À quel bureau appartient chacune de ces tâches~: enregistrer la correction du prix d'une transaction, approuver le P\&L du jour, choisir
l'option en espèces d'un dividende, exécuter un ordre~?
\end{exercise}

\begin{exercise}[$\star$]\label{exo:fm:operations:2}
Combien coûte un suspens dans les hypothèses du tutoriel, et quelle part de ce coût représente la pénalité~?
\end{exercise}

\begin{exercise}[$\star$]\label{exo:fm:operations:3}
Les écarts prenant 40 minutes en moyenne, quel est le plancher de la part en retard pour une fenêtre de 2.5 heures et pour une fenêtre de 10 heures~?
\end{exercise}

\begin{exercise}[$\star\star$]\label{exo:fm:operations:4}
Calculer la probabilité qu'un écart attende avec 11, 12 et 14 personnes ($\lambda=15$, $\mu=1.5$).
\end{exercise}

\begin{exercise}[$\star\star$]\label{exo:fm:operations:5}
Dans l'\cref{ex:fm:operations:call}, qu'estime devoir chaque partie, et pourquoi l'appel est-il contesté~?
\end{exercise}

\begin{exercise}[$\star\star$]\label{exo:fm:operations:6}
Pourquoi ajouter des personnes au-delà de treize réduit à peine les suspens sous T+1, et qu'est-ce qui les réduit~?
\end{exercise}

\begin{exercise}[$\star\star\star$]\label{exo:fm:operations:7}
\emph{Programmation.} Exécuter \texttt{fm\_ops.scenarios()}. Quelle modification du travail fait économiser le plus, diviser par deux le nombre d'écarts ou ramener
le travail à 30 minutes, et pourquoi le travail plus rapide demande-t-il davantage de personnes~?
\end{exercise}

\begin{exercise}[$\star\star\star$]\label{exo:fm:operations:8}
\emph{Trouver la faille.} «~Le passage à T+1 réduit de moitié notre fenêtre, nous allons donc doubler l'équipe des opérations.~»
\end{exercise}

\section{Problème~: un jour de moins}

\begin{problem}[{Problème du week-end --- un jour de moins}]\label{pb:fm:operations:1}
Un responsable des opérations doit présenter un plan pour le passage d'une firme d'un cycle de règlement-livraison de deux jours à un cycle d'un jour.

\medskip\noindent\textbf{Partie I --- L'organisation.}
\begin{enumerate}
\item Définir le front, le middle et le \omterm{def:fm:operations:offices}{back office}.
\item Pourquoi le \omterm{def:fm:operations:offices}{middle office} est-il indépendant du \omterm{def:fm:operations:offices}{front office}~?
\item Énoncer les règles américaines sur le cycle de règlement-livraison et sur l'affirmation le jour même.
\item Que constate le rapport de l'industrie après le passage américain~?
\end{enumerate}

\medskip\noindent\textbf{Partie II --- La file.}
\begin{enumerate}[resume]
\item Modéliser l'équipe des écarts comme une file d'attente et énoncer sa condition de stabilité.
\item Énoncer la \cref{prop:fm:operations:late} et démontrer le plancher.
\item Donner le plancher pour la fenêtre T+1 du tutoriel.
\item Donner la probabilité qu'un écart attende avec 11, 12 et 14 personnes.
\end{enumerate}

\medskip\noindent\textbf{Partie III --- Le dimensionnement.}
\begin{enumerate}[resume]
\item Donner l'effectif nécessaire pour un niveau de service de 95~\% sous chaque cycle.
\item Donner la part en retard, les suspens et le coût des suspens pour onze, douze et treize personnes sous T+1.
\item Quelle équipe minimise le coût total sous T+1, et quel est ce coût comparé à T+2~?
\item Donner le coût des deux modifications du travail présentées dans le \cref{tab:fm:operations:scenarios}.
\item Comment le profil de vieillissement évolue-t-il de onze à quatorze personnes~?
\item Qu'est-ce que le modèle laisse de côté, et dans quel sens cela le biaise-t-il~?
\end{enumerate}

\medskip\noindent\textbf{Partie IV --- Discipline et garanties.}
\begin{enumerate}[resume]
\item Définir un \omterm{def:fm:operations:discipline}{régime de discipline en matière de règlement} et donner les taux de pénalité de l'UE pour les actions et la dette souveraine.
\item Qu'est-ce qui a changé pour les buy-ins en 2023~?
\item Définir la \omterm{def:fm:operations:collateral}{gestion des garanties} (collateral management) et un \omterm{def:fm:operations:collateral}{litige sur les garanties}, puis résoudre l'\cref{ex:fm:operations:call}.
\item Définir un \omterm{def:fm:operations:kri}{indicateur clé de risque} et en donner quatre pour les opérations.
\item Énoncer le \emph{résultat nommé}~: l'effectif nécessaire pour maintenir 95~\% des écarts résolus avant le règlement sous un cycle de deux jours et sous un cycle d'un jour, et
le coût des suspens évité.
\item En deux phrases, rédiger le plan.
\end{enumerate}
\end{problem}

\section{Questions d'entretien}

\begin{interviewq}[$\star$ \iqroles{trader, développeur}]\label{iq:fm:operations:1}
Quelle est la différence entre le \omterm{def:fm:operations:offices}{middle office} et le \omterm{def:fm:operations:offices}{back office}~?
\end{interviewq}

\begin{interviewq}[$\star$ \iqroles{développeur}]\label{iq:fm:operations:2}
Qu'est-ce qu'un écart de rapprochement (break), et quelles en sont les causes habituelles~?
\end{interviewq}

\begin{interviewq}[$\star\star$ \iqroles{chercheur}]\label{iq:fm:operations:3}
Les écarts arrivent à raison de 15 par heure et prennent chacun 40 minutes. Combien de personnes vous faut-il au minimum, et pourquoi cela ne suffit-il pas~?
\end{interviewq}

\begin{interviewq}[$\star\star$ \iqroles{risque}]\label{iq:fm:operations:4}
Une contrepartie conteste votre appel de marge. Que faites-vous~?
\end{interviewq}

\begin{interviewq}[$\star\star$ \iqroles{trader, risque}]\label{iq:fm:operations:5}
Qu'a changé pour les opérations d'une société de trading le passage américain à T+1~?
\end{interviewq}

\begin{interviewq}[$\star\star\star$ \iqroles{développeur, chercheur}]\label{iq:fm:operations:6}
Concevez un système qui réduit les suspens de règlement-livraison sous un cycle d'un jour.
\end{interviewq}
